\color{unidarkgreen}{\section[Порядок технического обслуживания изделия ]{Порядок технического обслуживания изделия }\label{app:pto}}
\color{black}

\par
В таблице \ref{apppto:tbl1} приведены виды работ при соответствующих проверках.


\setlength{\extrarowheight}{0.06cm}
\small
\begin{longtable}{
  |>{\raggedright\arraybackslash}m{2cm}
  |>{\raggedright\arraybackslash}m{14.8cm}|
}
\caption{Виды работ при проверке устройства\hfill\vspace{-0.5\baselineskip}}\label{apppto:tbl1}\\ 
\hline
\rowcolor{gray!30}
Виды проверок & Виды работ при проверке \\
\hline
\endfirsthead

\caption*{\hspace{3pt}\emph{Продолжение таблицы \ref{apppto:tbl1}\hfill\vspace{-0.5\baselineskip}}} \\
\hline
\rowcolor{gray!30}
Виды проверок & Виды работ при проверке \\ 
\hline
\endhead
\hline
\endfoot
\endlastfoot

Н, К1, В, К & а) внешний осмотр: отсутствие внешних следов ударов, потеков воды, в том числе высохших, отсутствие налета окислов на металлических поверхностях, отсутствие запыленности, осмотр рядов зажимов входных и выходных сигналов, разъемов интерфейса связи в части состояния их контактных поверхностей, осмотр элементов управления на отсутствие механических повреждений \\ \hline
		  В & б) внутренний осмотр: чистка от пыли; осмотр элементов цепей с точки зрения наличия следов перегревов, ослабления паяных соединений из-за появления трещин, наличия окисления; контроль сочленения разъемов и механического крепления элементов, затяжка винтовых соединений \\ \hline
Н, К1, В, К & в) измерение сопротивления изоляции независимых цепей (кроме цепей интерфейса связи) по отношению к корпусу и между собой \\ \hline
	   Н, В & г) испытания электрической прочности изоляции независимых цепей (кроме цепей интерфейса связи) по отношению к корпусу и между собой \\ \hline
          Н & д) проверка работоспособности дискретных входов, выходных реле и светодиодов терминала \\ \hline
   Н, К1, В & е) задание (или проверка) требуемой конфигурации устройства в соответствии с принятыми проектными решениями и техническими характеристиками (функциями) устройства \\ \hline
   Н, К1, В & ж) задание (или проверка) уставок устройства в соответствии с заданной конфигурацией \\ \hline
   Н, К1, В & з) проверка правильности отображения значений и фазовых углов токов (напряжений), поданных от постороннего источника \\ \hline
   Н, К1, В & и) проверка параметров (уставок) срабатывания и коэффициентов возврата каждого ИО при подаче на входы устройства тока (напряжения) от постороннего источника; контроль состояния светодиодов при срабатывании \\ \hline
          Н & к) проверка срабатывания устройства на рабочих уставках и определение изменения параметров срабатывания при напряжении оперативного тока, равном 0,8 и 1,1 Uном \\ \hline
   Н, К1, В & л) проверка времени срабатывания защиты и автоматики на соответствие заданным уставкам по времени \\ \hline
          Н & м) проверка отсутствия ложных действий при снятии и подаче напряжения оперативного тока с повторным включением через 0,5 с при минимальном значении диапазона уставок с подачей тока (напряжения), равного 0,8 тока (напряжения) срабатывания \\ \hline
	   Н, В & н) проверка взаимодействия ИО и логических цепей защиты с контролем состояния всех контактов выходных реле и визуальным контролем состояния светодиодов и ламп сигнализации. Проверка проводится при напряжении питания оперативного тока, равном 0,8 Uном, и создании условий для поочередного срабатывания каждого ИО и подачи необходимых сигналов на дискретные входы защиты \\ \hline
Н, К1, В, К & о) проверка управляющих функций защиты с воздействием контактов выходного реле в цепи управления коммутационным аппаратом \\ \hline
	   Н, В & п) проверка функций регистрации событий, осциллографирования сигналов, отображения параметров защиты \\ \hline
Н, К1, В, К & р) проверка управления коммутационным аппаратом присоединения (включить/отключить) \\ \hline
   Н, К1, В & с) проверка взаимодействия с другими устройствами защиты, электроавтоматики, управления и сигнализации с воздействием на коммутационный аппарат \\ \hline
Н, К1, В, К & т) проверка рабочим током \\ \hline
	   
\end{longtable}
\normalsize
\setlength{\extrarowheight}{0.0cm}

\par
\textbf{Профилактический контроль}

\par
Терминалы имеют встроенную систему самодиагностики и не требуют периодического тестирования.
Самодиагностика обеспечивает локализацию повреждения с точностью до блока терминала.

\par
Особое внимание при проведении профилактического контроля следует уделить проверке винтов на клеммах терминала и на рядах зажимов шкафа.

\par
При проведении работ по профилактическому контролю рекомендуется измерить переменные токи и напряжения, подводимые к зажимам шкафа, и провести сравнение их с показаниями токов и напряжений на дисплее терминала.
При соответствии показаний дальнейшую проверку уставок защит допускается не проводить.

\par
При проведении работ по профилактическому контролю следует проверить исправность дискретных входов терминала, а также замыкание контактов выходных реле шкафа.
Перед выполнением проверки необходимо принять меры для исключения действия шкафа во внешние цепи.

\par
Проверку исправности дискретных входов, выведенных на клеммы шкафа, а также переключателей и кнопок шкафа следует проводить в соответствие с рекомендациями изготовителя терминала.
При этом на дискретный вход необходимо подать напряжение, соответствующее логической «1», установить переключатель в необходимое положение и/или нажать кнопку, контролируя на дисплее терминала в соответствующем меню тестового режима состояние контролируемого входа.

\par
Проверку исправности дискретных выходов проводят по рекомендациям изготовителя терминала, при этом логически имитируется срабатывание внутреннего сигнала терминала, запрограммированное на выходное реле.

\par
\textbf{Профилактическое восстановление}

\par
Обслуживающий шкаф персонал может самостоятельно провести ремонт или замену внешних реле шкафа, переключателей, светосигнальной арматуры и т.д.

\par
\textbf{ВНИМАНИЕ!} В случае обнаружения дефектов в терминале или в устройстве связи с ПК, необходимо немедленно поставить в известность предприятие-изготовитель.
Восстановление вышеуказанной аппаратуры может производить только специально подготовленный персонал.

\clearpage